\section{CONCLUSÕES}

As conclusões seguem a ordem das perguntas de pesquisa. Todas são associações e valem com as ressalvas da Seção 8.

A medição do texto é utilizável, com restrições. Em 297 sentenças rotuladas por um pesquisador, a acurácia populacional do sentimento é de 87\% no modelo A e de 85\% no B. O kappa ponderado populacional é de 0,66 e 0,51. Pela regra fixada antes da execução do modelo B, com diferença de kappa de $-0{,}141$ (IC 95\% de $-0{,}272$ a $-0{,}012$), o modelo principal é A. O modelo B acompanha todos os resultados. Os dois modelos ordenam sentenças e anúncios de forma semelhante, mas não são intercambiáveis. O sentimento por aspecto concorda de forma moderada e positiva com as subnotas oficiais. Apenas limpeza, localização e custo-benefício atendem ao critério de confiança. Check-in, comunicação, ruído e conforto, segurança, precisão do anúncio, comodidades e satisfação geral pedem cautela. Um segundo método de tópicos (LDA) reproduz o sinal nos aspectos que consegue checar. Entre dois anotadores, a concordância é alta no sentimento (kappa de 0,86) e moderada no aspecto (0,62).

O texto das avaliações revela o que se associa à nota. Nos dois modelos, comodidades é o aspecto de maior parcela do $R^2$ da nota global (42\% em A e 48\% em B), seguida de limpeza (20\% em A e 19\% em B). Isso se mantém em todos os segmentos e sensibilidades. Como comodidades pede cautela, a leitura mais sustentada é a dos tópicos finos. Treze das 15 queixas de maior peso coincidem nos dois modelos (banheiro e chuveiro, ar-condicionado, sujeira, cozinha, cheiro, cama, \textit{wi-fi} e resposta do anfitrião) e se associam a quedas de 10 a 26 pontos no percentil da nota.

A decomposição é aplicável ao anúncio. O diagnóstico por anúncio compara cada aspecto com os anúncios do mesmo bairro e tipo, mostra a evidência de cada aspecto e as queixas e elogios recorrentes e indica onde investigar, sem prever a nota nem prescrever causas.

Um índice único não acrescenta poder explicativo. O índice sintético reordena os anúncios em relação à nota oficial e discrimina os anúncios de nota 5,0 que o possuem (5.142 de 9.373), mas não acrescenta poder explicativo sobre os índices de volume de avaliações além da nota oficial, em nenhum dos dois modelos. O resultado é nulo e secundário, com duas leituras que não se excluem: o sentimento tem pouca razão teórica para prever uma contagem de avaliações; o valor do texto está na decomposição por aspecto. Os coeficientes do índice dentro da amostra não são robustos e ficam fora desta síntese.

O estudo não conclui que o índice sintético seja uma nota melhor que a oficial, que corrigir um problema apontado em uma queixa eleve a nota, nem que exista efeito sobre receita ou ocupação.

Alguns desdobramentos seguem diretamente das limitações. Desfechos econômicos observados, como calendário de disponibilidade e preços praticados ao longo do tempo, permitiriam testar o índice sintético com um desfecho que não seja índice de volume. Um painel temporal permitiria acompanhar o mesmo anúncio e aproximar-se de um contraste dentro do anúncio. A rotulagem por mais anotadores, em especial dos aspectos que pedem cautela, reduziria o erro de referência da atribuição de aspecto. A replicação em outras cidades testaria a generalidade da decomposição.
